% Based on templates/lecture.tex; actual class 11, 2026-10-28 (Week 11).
% Sources: Lec4_week11.pdf pp. 4--51; Lec4_recorded.pdf pp. 4--43.
% Transcript: 81346 - 0_id8644hz - PID 117.txt (no timestamps).
% Live pp. 24--25 duplicate; pp. 37 and 44 are Optional.
% Source scope and numerical examples verified; no dialogue checks retained.

\begingroup
\linespread{1.15}\selectfont
\setlength{\parskip}{0.4em}
\renewcommand{\recordid}[1]{}
\renewcommand{\relatedexercise}[2]{\hyperref[exercise:#1]{#2}}
\section{第 11 节课：Attention 与 Transformers}
\label{lecture:SC4002-L11}

\noteintro{
\textbf{日期与来源：}2026-10-28（Week 11）；\texttt{Lec4\_week11.pdf} pp.~4--51、\texttt{Lec4\_recorded.pdf} pp.~4--43；字幕 81346（无时间戳）。重复内容合并，Optional 内容随节标明；笔记待审阅。
}

\textbf{本讲主线。}从固定向量概括源句，转向按需读取各个位置；再让词与词交换信息，组合成 Transformer，最后按任务确定可见范围与输出方式。

\lecturetopic{SC4002-L11-T01}{Cross-attention：每一步按需读取源句}
\nopagebreak[4]
\textbf{来源：}现场版 pp.~4--17；录播版 pp.~4--21。
\nopagebreak[4]

将 \texttt{red car} 翻译为 \texttt{voiture rouge} 时，两个生成步骤需要关注不同源词。只把源句压入 encoder 最后一个状态，会形成 information bottleneck（信息瓶颈）。Attention 保留整组 encoder 状态 $\mathbf h_1,\ldots,\mathbf h_n$，让每个 decoder 步骤重新读取所需信息。

\begin{center}
\small
\begin{tabular}{@{}lll@{}}
\toprule
角色 & 本例中的对象 & 用途 \\
\midrule
Query（查询） & 当前 decoder 状态 $\mathbf s_t$ & 表达这一步的信息需求 \\
Key（键） & 各 encoder 状态 $\mathbf h_i$ & 与 query 计算匹配分数 \\
Value（值） & 各 encoder 状态 $\mathbf h_i$ & 提供加权汇总的内容 \\
\bottomrule
\end{tabular}
\end{center}

本例中 key 与 value 共用 $\mathbf h_i$；两者承担不同功能，不要求一般模型也使用相同向量。对同维列向量，未缩放的点积 attention 为
\begin{align*}
 e_{t,i}&=\mathbf s_t^\top\mathbf h_i, &
 \alpha_{t,i}&=\frac{\exp(e_{t,i})}{\sum_{j=1}^{n}\exp(e_{t,j})},\\
 \mathbf c_t&=\sum_{i=1}^{n}\alpha_{t,i}\mathbf h_i, &
 \mathbf p_t&=\operatorname{softmax}\!\left(W_o[\mathbf s_t;\mathbf c_t]+\mathbf b_o\right).
\end{align*}
其中 $\mathbf c_t$ 是 context vector（上下文向量），对应录播符号 $\mathbf a_t$；$[\,;\,]$ 表示拼接。第一个 softmax 在\textbf{源位置}上分配权重，最后一个 softmax 在\textbf{目标词表}上给出预测概率。

\begin{center}
% Original diagram based on Lec4_recorded.pdf pp. 7, 14, 20.
\begin{tikzpicture}[
 >=Latex,font=\small,
 box/.style={draw=noteBlue,rounded corners,fill=blue!3,minimum height=.65cm,align=center},
 flow/.style={->,thick,noteBlue}
]
 \node[box,minimum width=2.1cm] (q) at (0,0) {Query $\mathbf q$};
 \node[box,minimum width=2.1cm] (k) at (0,1.1) {Keys $\{\mathbf k_i\}$};
 \node[box,minimum width=2.2cm] (score) at (3,0) {匹配分数 $e_i$};
 \node[box,minimum width=2.2cm] (alpha) at (6.3,0) {权重 $\alpha_i$};
 \node[box,minimum width=2.7cm] (out) at (10.3,0) {$\mathbf c=\sum_i\alpha_i\mathbf v_i$};
 \node[box,minimum width=2.7cm] (v) at (10.3,1.1) {Values $\{\mathbf v_i\}$};
 \draw[flow] (q) -- (score);
 \draw[flow] (k.east) -| (score.north);
 \draw[flow] (score) -- node[above,font=\footnotesize] {softmax} (alpha);
 \draw[flow] (alpha) -- node[above,font=\footnotesize] {加权} (out);
 \draw[flow] (v) -- (out);
\end{tikzpicture}

\end{center}
{\small\textit{原创流程图，依据录播版 pp.~7、14、20。Key 用于匹配，value 用于汇总；这里一个 query 产生一个上下文向量。}}

注意力权重随 query 改变，不必只选择一个源词；它可以融合多处证据。把各生成步骤的权重排成矩阵，可观察 soft alignment（软对齐），但不能仅凭权重证明完整的因果解释。直接读取各源位置绕过了固定摘要瓶颈，也为远处状态提供较短的梯度通路；这有助于缓解梯度消失，不保证所有优化问题都消失。

\textbf{讲义例题：}\relatedexercise{SC4002-L11-E01}{\texttt{red car} 的两步读取与输出层}。

\lecturetopic{SC4002-L11-T02}{Self-attention：为每个词形成上下文表示}
\nopagebreak[4]
\textbf{来源：}现场版 pp.~18--26；录播版 pp.~22--34。

Self-attention（自注意力）让同一序列的各位置交换信息。例如 \texttt{not very good} 中，good 的表示需要结合前面的修饰与否定。它的直接输出是\textbf{每个位置的新表示}，还不是整句类别或下一词。

\subsubsection{投影、匹配与汇总}

批量公式采用一词一行的约定。设 $X\in\mathbb R^{n\times d}$，则
\[
 Q=XW_Q,\quad K=XW_K,\quad V=XW_V,
 \qquad
 A=\operatorname{softmax}_{\rm row}\!\left(\frac{QK^\top}{\sqrt{d_k}}\right),
 \qquad O=AV.
\]
$W_Q,W_K,W_V$ 是可学习参数，$Q,K,V$ 是当前输入产生的表示。三者来自同一序列，但投影不同，数值无需相同。$\operatorname{softmax}_{\rm row}$ 表示每个 query 分别在 key 位置上归一化。缩放因子 $\sqrt{d_k}$ 控制点积尺度；现场版算例使用缩放，录播的基础算例没有使用，不能混用。

\begin{center}
\small
\begin{tabular}{@{}lll@{}}
\toprule
对象 & 形状 & 含义 \\
\midrule
$W_Q,W_K$ & $d\times d_k$ & 查询与键的投影参数 \\
$W_V$ & $d\times d_v$ & 内容投影参数 \\
$Q,K$；$V$ & $n\times d_k$；$n\times d_v$ & 每个位置的三种表示 \\
$A$ & $n\times n$ & 第 $i$ 行查询、第 $j$ 列被读取 \\
$O$ & $n\times d_v$ & 每个位置一个 $d_v$ 维输出 \\
\bottomrule
\end{tabular}
\end{center}

单个位置的输出为 $\mathbf o_i=\sum_j A_{ij}\mathbf v_j$：加权汇总发生在\textbf{被读取的位置}上，保留 value 的特征维度。对全部 query 重复该运算，便得到 $n$ 行输出；整句汇聚需要后续 pooling，不能把两者混为一谈。

\subsubsection{与 RNN 的区别及边界}

RNN 的当前状态依赖前一状态；self-attention 同一层的各 query 都读取上一层已知表示，因此可并行计算。全连接 attention 使任意两个位置能在一层中直接交互，讲义的 $O(1)$ 指\emph{交互路径长度}。全体配对仍形成 $n\times n$ 的分数矩阵，总计算量不是 $O(1)$。因果生成还要限制可见范围，见\hyperref[topic:SC4002-L11-T04]{任务与掩码}。

\textbf{讲义例题：}\relatedexercise{SC4002-L11-E02}{\texttt{not very good} 的缩放 attention}；\relatedexercise{SC4002-L11-E03}{录播三维例题的精确 softmax}。

\lecturetopic{SC4002-L11-T03}{Transformer block：交换、加工与逐层更新}
\nopagebreak[4]
\textbf{来源：}现场版 pp.~28--41；录播版 p.~36。多头组合公式在现场版 p.~37 标为 Optional。

\subsubsection{位置与逐位置非线性}

完整 self-attention 本身不内置词序：交换输入位置会相应交换输出，无法仅凭这一步区分 \texttt{dog bites man} 与 \texttt{man bites dog} 的顺序关系。讲义将同维 position embedding（位置嵌入）加到词向量：
\[
 \widetilde{\mathbf x}_i=\mathbf x_i+\mathbf p_i.
\]
Learned absolute position embeddings（可学习的绝对位置嵌入）为每个位置训练一个向量，灵活但受位置表长度限制；表外位置没有可直接查用的向量。本讲不展开其他位置编码公式。

Attention 汇聚跨位置内容后，以共享的 feed-forward network（前馈网络，FFN）加工每个位置。以下单个位置公式采用列向量：
\[
 \operatorname{FFN}(\mathbf o_i)
 =W_2\operatorname{ReLU}(W_1\mathbf o_i+\mathbf b_1)+\mathbf b_2.
\]
Attention 负责位置间的信息交换，FFN 负责逐位置的特征变换。Softmax 本身已有非线性；FFN 补充的是对表示的显式非线性加工，不应把整个 attention 称为线性运算。

\subsubsection{多头、残差与归一化}

Multi-head attention（多头注意力）让不同头使用各自的投影，在不同表示空间中计算权重，然后拼接并映射回模型维度：
\[
 \mathrm{head}_r=\operatorname{Attention}(Q_r,K_r,V_r),\qquad
 \operatorname{MultiHead}=\operatorname{Concat}(\mathrm{head}_1,\ldots,\mathrm{head}_H)W^O.
\]
讲义以实体关联、语法关联和代词指代说明不同头可能关注的内容；这些是可能学到的模式，不是固定分工，也不保证每个头都有独立可解释的功能。

Residual connection（残差连接）保留输入的直接通路；layer normalization（层归一化）在每个位置的特征维上归一化，帮助稳定训练。按讲义图示的 Add \& Norm 顺序，一个 encoder block 可写为
\begin{align*}
 U&=\operatorname{LayerNorm}\!\left(X+\operatorname{MultiHead}(X)\right),\\
 Y&=\operatorname{LayerNorm}\!\left(U+\operatorname{FFN}(U)\right).
\end{align*}
这里 $X$ 是该层输入；第一层输入包含词与位置表示。先交换信息，再加工特征，每个子层都有 Add \& Norm；继续堆叠多个 block，便逐层更新上下文表示。

\lecturetopic{SC4002-L11-T04}{任务决定可见范围与输出方式}
\nopagebreak[4]
\textbf{来源：}现场版 pp.~42--50（p.~44 为 Optional）；录播版 pp.~39--41。

\begin{center}
\small
\begin{tabularx}{\textwidth}{@{}lXX@{}}
\toprule
架构 & 信息来源 & 输出与讲义中的用途 \\
\midrule
Encoder-only & 完整输入的双向上下文 & 整句池化后分类；逐词分类用于 NER；BERT-style \\
Decoder-only & 当前及之前的位置 & 逐位置预测下一词；GPT-style 续写 \\
Encoder--decoder & 完整源句与目标前缀 & 翻译、摘要等条件生成；T5-style \\
\bottomrule
\end{tabularx}
\end{center}

\subsubsection{完整输入与因果掩码}

对 \texttt{not very good} 做情感分类时，整句已知：encoder 产生逐词表示，经 pooling（汇聚）得到句向量，再接类别输出层。训练目标可用类别交叉熵。

预测下一词时，位置 $i$ 不得读取未来位置 $j>i$。Causal mask（因果掩码）在 softmax \emph{之前}作用于分数：
\[
 A=\operatorname{softmax}_{\rm row}\!\left(\frac{QK^\top}{\sqrt{d_k}}+M\right),
 \qquad
 M_{ij}=\begin{cases}0,&j\le i,\\-\infty,&j>i.\end{cases}
\]
被屏蔽位置的 softmax 后权重为 $0$；其余位置重新归一化。现场版 p.~48 的图混合展示了屏蔽值与概率，复习时必须区分这两个阶段。

最终逐位置表示 $\mathbf h_t$ 经输出层得到词表分布：
\[
 \mathbf z_t=W_{\rm vocab}\mathbf h_t+\mathbf b,\qquad
 \mathbf p_t=\operatorname{softmax}(\mathbf z_t),\qquad
 L=-\sum_t\log p_{t,y_t}.
\]
$y_t$ 表示该位置应预测的真实下一词。完整训练序列提供标签，因果掩码允许并行计算多个训练位置而不泄漏未来；生成时后续词尚未产生，仍需逐词选择并接续。Softmax 输出是概率；若需要 log-probabilities，须再取对数或使用 log-softmax。

\textbf{讲义例题：}\relatedexercise{SC4002-L11-E04}{同一输入在 full 与 causal attention 下的差异}。

\subsubsection{Encoder--decoder 的完整信息流}

Encoder 通过完整 self-attention 得到整组源句表示。每个 decoder block 先用 masked self-attention 处理目标前缀，再用 cross-attention 查阅源句，最后经过 FFN；各子层均带 Add \& Norm。

\begin{center}
% Original diagram based on Lec4_week11.pdf pp. 41--44.
% One encoder and one decoder block shown; Add & Norm follows each sublayer.
\begin{tikzpicture}[
 >=Latex,font=\small,
 box/.style={draw=noteBlue,fill=blue!3,rounded corners,minimum width=4.2cm,minimum height=.65cm,align=center},
 flow/.style={->,thick,noteBlue}
]
 \node (src) at (0,0) {源句：词向量 $+$ 位置};
 \node[box] (encsa) at (0,1.2) {Full self-attention};
 \node[box] (encff) at (0,2.4) {FFN};
 \node[box] (encout) at (0,3.6) {Encoder 最终输出};
 \node (tgt) at (8,0) {目标前缀：词向量 $+$ 位置};
 \node[box] (decsa) at (8,1.2) {Masked self-attention};
 \node[box,fill=ntuRed!5,draw=ntuRed] (cross) at (8,2.4) {Cross-attention};
 \node[box] (decff) at (8,3.6) {FFN};
 \node[box] (pred) at (8,4.8) {词表投影 $+$ softmax};
 \draw[flow] (src) -- (encsa);
 \draw[flow] (encsa) -- (encff);
 \draw[flow] (encff) -- node[left,font=\footnotesize] {可堆叠} (encout);
 \draw[flow] (tgt) -- (decsa);
 \draw[flow] (decsa) -- node[right,font=\footnotesize] {$Q$} (cross);
 \draw[flow] (cross) -- (decff);
 \draw[flow] (decff) -- node[right,font=\footnotesize] {可堆叠} (pred);
 \draw[->,thick,ntuRed] (encout.east) -- (3.5,3.6)
   -- node[right,font=\footnotesize] {$K,V$} (3.5,2.4) -- (cross.west);
 \node[font=\footnotesize,align=center,text width=4.5cm] at (0,4.85)
   {每个 attention / FFN 子层后\\均有 Add \& Norm};
\end{tikzpicture}

\end{center}
{\small\textit{原创结构图，依据现场版 pp.~41--44。实线表示前向计算；Add \& Norm 在每个子层后执行，图中合并标示以突出信息来源。}}

Cross-attention 的 query 来自 decoder 当前表示，keys 与 values 来自 encoder 最终输出；encoder 输出可供多个 decoder 层反复读取。Masked self-attention 与 cross-attention 在一次前向计算中依次执行，整网按目标序列损失联合训练，无须先单独训练其中一个模块。

录播末尾以历史实验展示翻译的 BLEU、文档生成的 perplexity 与 ROUGE-L，以及 GLUE 语言理解结果。它们说明应用价值，不构成所有任务和计算预算下 Transformer 必然最优的结论。

\subsection*{一分钟回顾}

Cross-attention 让每个生成步骤按需读取源句；self-attention 为每个输入位置融合上下文。位置嵌入补充顺序，FFN 加工特征，多头提供多种读取方式，Add \& Norm 支持多层训练。理解任务可读完整输入，生成任务须限制未来可见性；attention 输出还需通过任务输出层，才能得到类别或词表概率。
\endgroup
